\begin{frame}{왜 BC만으로는 부족한가: 복합 오차의 눈덩이}
  \begin{itemize}
    \item 학습은 \textbf{전문가가 실제로 방문한 상태}들에서만 이뤄진다.
    \item 배포된 정책은 작은 실수 하나로 \textbf{전문가가 가본 적 없는}
      상태로 빗나갈 수 있다.
    \item 낯선 상태에선 정책을 \textbf{배운 적 없으므로} 더 큰 실수.
    \item 그 실수가 또 다른 낯선 상태로 이어진다.
  \end{itemize}
  \vskip0.6em
  \begin{block}{\kb{복합 오차}(compounding error)}
    오차가 시간이 지날수록 눈덩이처럼 불어나는 문제.
  \end{block}
  \vskip0.5em
  \kq{어긋남은 한 스텝의 실수($\pi\ne\pi^E$ 인 스텝)가 경로를 시연 밖으로
      빼내는 순간 시작된다.}
\end{frame}
